\chapter{Introduction}
\label{ch:intro}

Static code analysis has been a part of software engineering practice since the 1970s, when Johnson introduced \texttt{lint} for checking C programs~\cite{johnson1977lint}. The technique examines source code without executing it, identifying potential defects through analysis of code structure, control flow, and data flow. Over the decades, static analysis has evolved from simple pattern matching to sophisticated tools capable of inter-procedural analysis, and has become a standard component of industrial software development processes~\cite{bessey2010coverity}.

CodeChecker~\cite{codechecker} is an open-source static analysis infrastructure built on the LLVM/Clang Static Analyzer~\cite{clangsa} toolchain. It provides a platform for managing static analysis results from multiple analyzers, including Clang Static Analyzer, Clang Tidy, Cppcheck, and various other tools supporting C/C++, Java, Python, and JavaScript. CodeChecker offers both command-line interfaces and a web-based visualization tool, enabling teams to store, query, filter, and compare analysis results across multiple code revisions and analysis runs.

While static analysis tools identify potential code defects, they operate independently of test execution. Test coverage analysis~\cite{yang2009coverage} measures which lines of source code are executed during test runs, providing developers with insights into how thoroughly their test suites exercise the codebase. Tools such as \texttt{gcov}~\cite{gcov} and \texttt{lcov}~\cite{lcov} are widely used in the C/C++ ecosystem for this purpose. The integration of code coverage visualization with static analysis results presents an opportunity to correlate static analysis findings with actual test coverage, enabling developers to identify areas where both static analysis warnings exist and test coverage is insufficient.

This thesis presents the design and implementation of a code coverage visualization feature integrated into CodeChecker. The feature allows users to convert existing test coverage data (in LCOV format) into CodeChecker's internal format using the \texttt{report-converter} tool, store it alongside static analysis results, and visualize coverage information through an interactive web interface that highlights covered and uncovered lines of source code.

The implementation required expertise in static analysis toolchains, database schema design with SQLAlchemy~\cite{sqlalchemy}, web application development with Vue.js~\cite{vuejs} and CodeMirror~\cite{codemirror}, API design with Apache Thrift~\cite{thrift}, and coverage instrumentation techniques. The feature was validated on the Apache Xerces-C++ XML parser~\cite{xerces}, a real-world C++ project with 80 test cases and over 580 source files.

The remainder of this thesis is structured as follows. Chapter~\ref{ch:user} provides user documentation describing the coverage workflow and the web interface. Chapter~\ref{ch:impl} presents the technical implementation details, including the database schema, API endpoints, LCOV converter, and frontend components. Chapter~\ref{ch:sum} summarizes the contributions and discusses potential future enhancements.

\section{Related Work}

Several tools and platforms address code coverage visualization, each with different integration approaches.

\textbf{LCOV genhtml}~\cite{lcov} generates static HTML reports from LCOV \texttt{.info} files, producing a directory-browsable view with line and function coverage percentages. It is the de facto standard for C/C++ coverage visualization. Our implementation draws inspiration from its visual design (color-coded coverage bars, directory drill-down) but integrates the data into CodeChecker's existing web application rather than producing standalone HTML files.

\textbf{Codecov} and \textbf{Coveralls} are cloud-hosted services that accept coverage data from CI pipelines and provide web dashboards with coverage trends, pull request annotations, and team-level reporting. They support many languages and coverage formats but operate as external services, requiring data to leave the development environment. Our approach keeps all data within the CodeChecker server, which is important for organizations with strict data governance requirements.

\textbf{SonarQube} integrates coverage visualization with static analysis in a single platform, similar to our goal. However, SonarQube uses its own analysis engine and is a separate product from the compiler toolchain. CodeChecker, by contrast, is built on LLVM/Clang and integrates directly with the compiler's static analyzers, making the addition of compiler-generated coverage data a natural extension.

\textbf{gcov}~\cite{gcov} is the GCC coverage tool that produces per-file \texttt{.gcov} annotated source listings. CodeChecker's \texttt{report-converter} already included a \texttt{gcov} converter before this work. Our LCOV converter complements it by accepting the more widely used LCOV format, which aggregates data from multiple \texttt{gcov} files and is the standard output of CI coverage pipelines.

The key differentiator of our approach is the integration of coverage data into an existing static analysis platform through the \texttt{report-converter} pattern, allowing developers to view both static analysis warnings and test coverage in the same interface without adopting a new tool.
